% Prepared for Optimization Letters using the publisher's December 2024 template.
% The publisher class is supplied unchanged. This is a single-file manuscript.
\documentclass[pdflatex,sn-mathphys-num]{sn-jnl}
\usepackage[T1]{fontenc}
\usepackage{lmodern}
\usepackage{amsmath,amssymb,amsthm}
\usepackage{booktabs,array}
\usepackage{graphicx}
\usepackage{xcolor}
\usepackage{microtype}
% BEGIN COMPACT LAYOUT v13
% Keep the publisher class, body fonts, line spacing and page margins unchanged.
\usepackage{etoolbox}
\makeatletter
\apptocmd{\normalsize}{%
  \abovedisplayskip=7pt plus1pt minus1pt
  \belowdisplayskip=7pt plus1pt minus1pt
  \abovedisplayshortskip=2pt plus1pt
  \belowdisplayshortskip=3pt plus1pt minus1pt
}{}{\PackageError{compact-layout}{Normal-size spacing patch failed}{}}
\apptocmd{\th@thmstyleone}{%
  \thm@preskip=8pt plus1pt minus1pt
  \thm@postskip=8pt plus1pt minus1pt
}{}{\PackageError{compact-layout}{Theorem spacing patch failed}{}}
\apptocmd{\th@thmstyletwo}{%
  \thm@preskip=8pt plus1pt minus1pt
  \thm@postskip=8pt plus1pt minus1pt
}{}{\PackageError{compact-layout}{Remark spacing patch failed}{}}
\renewcommand\section{\@startsection{section}{1}{\z@}%
  {-10pt plus-2pt minus-1pt}{5pt}{\sectionfont}}
\renewcommand\subsection{\@startsection{subsection}{2}{\z@}%
  {-8pt plus-2pt minus-1pt}{4pt}{\subsectionfont}}
\renewcommand\subsubsection{\@startsection{subsubsection}{3}{\z@}%
  {-8pt plus-2pt minus-1pt}{4pt}{\subsubsectionfont}}
\patchcmd{\@maketitle}{\vskip21pt}{\vskip8pt}{}{\PackageError{compact-layout}{Title top patch failed}{}}
\patchcmd{\@maketitle}{\vskip20pt}{\vskip12pt}{}{\PackageError{compact-layout}{Title gap patch failed}{}}
\patchcmd{\@maketitle}{\vskip20pt}{\vskip12pt}{}{\PackageError{compact-layout}{Author gap patch failed}{}}
\patchcmd{\@maketitle}{\vskip24pt}{\vskip12pt}{}{\PackageError{compact-layout}{Email gap patch failed}{}}
\patchcmd{\@maketitle}{\vskip24pt}{\vskip12pt}{}{\PackageError{compact-layout}{Abstract gap patch failed}{}}
\patchcmd{\@maketitle}{\vskip36pt}{\vskip10pt}{}{\PackageError{compact-layout}{Title end patch failed}{}}
\AtBeginDocument{%
  \normalsize
}
\makeatother
\setlength{\floatsep}{9pt plus2pt minus1pt}
\setlength{\textfloatsep}{9pt plus2pt minus1pt}
\setlength{\intextsep}{9pt plus2pt minus1pt}
\setlength{\belowcaptionskip}{4pt}
\setlength{\bibsep}{4pt}
% END COMPACT LAYOUT v13
\hypersetup{hidelinks,pdftitle={A Fourth-Power Relative-Error Threshold for Two-Periodic Heavy Ball},
pdfauthor={Bahadir Yuzbasi and Mohammad Arashi},pdfsubject={Optimization Letters manuscript}}
\theoremstyle{thmstyleone}
\newtheorem{theorem}{Theorem}
\newtheorem{lemma}{Lemma}
\newtheorem{corollary}{Corollary}
\theoremstyle{thmstyletwo}
\newtheorem{remark}{Remark}
\raggedbottom
\newcommand{\R}{\mathbb R}
\newcommand{\spec}{\operatorname{spec}}
\newcommand{\tr}{\operatorname{tr}}
\newcommand{\diag}{\operatorname{diag}}
\newcommand{\crit}{\mathrm{crit}}
\newcommand{\cT}{\mathsf T}
\newcommand{\cK}{\mathsf K}
\newcommand{\cG}{\mathsf G}
\newcommand{\norm}[1]{\left\lVert #1\right\rVert}

\begin{document}
\title[Relative errors in two-periodic heavy ball]{A Fourth-Power Relative-Error Threshold for Two-Periodic Heavy Ball}
\author*[1,2]{\fnm{Bahad{\i}r} \sur{Y\"uzba\c{s}{\i}}}\email{bahadir.yuzbasi@inonu.edu.tr}
\author[3]{\fnm{Mohammad} \sur{Arashi}}\email{arashi@um.ac.ir}
\affil*[1]{\orgdiv{Department of Econometrics}, \orgname{\.{I}n\"on\"u University},
\orgaddress{\city{Malatya}, \country{T\"urkiye}}}
\affil[2]{\orgdiv{Department of Statistics}, \orgname{F{\i}rat University},
\orgaddress{\postcode{23199}, \city{Elaz{\i}\u{g}}, \country{T\"urkiye}}}
\affil[3]{\orgdiv{Department of Statistics, Faculty of Mathematical Sciences},
\orgname{Ferdowsi University of Mashhad},
\orgaddress{\city{Mashhad}, \country{Iran}}}
\abstract{\unboldmath
Periodic heavy-ball coefficients can exploit a gap in the Hessian
spectrum of a strongly convex quadratic. We study how accurately the
gradient must be computed to preserve this advantage. For two equal-width
spectral bands and a fixed condition number $\kappa>1$, we consider the
established two-periodic coefficients tuned for exact gradients, with
zero initial momentum. We prove that the supremum of relative-error levels
preserving a per-query rate below
$q=(\sqrt\kappa-1)/(\sqrt\kappa+1)$ has order $R^4$ as the normalized
spectral gap $R$ tends to zero. Both bounds have explicit constants.
The lower bound is uniform in dimension and allows history-dependent
vector errors that mix Hessian eigenspaces. Its proof uses a regularized
quadratic energy and gives an estimate for every iterate. A scalar
multiplicative error yields the matching upper order for the same
initialization. Controlled two-band examples illustrate the scalar fourth-power
scaling and cyclical heavy ball's smaller terminal errors under a fixed query budget.
On a least-squares example built from observed stock index returns,
cyclical heavy ball achieves smaller terminal distances and objective gaps than stationary heavy ball, gradient
descent and quadratically tuned Nesterov, after an early overshoot.
The theorem concerns a benchmark-rate
threshold for fixed coefficients, rather than a stability threshold
or the full rate--robustness tradeoff under inexactness.}
\keywords{heavy-ball method, cyclical step sizes, inexact gradients, spectral gaps, robustness, quadratic optimization}
\maketitle
\noindent\textit{ORCID:} Bahad{\i}r Y\"uzba\c{s}{\i}: \url{https://orcid.org/0000-0002-6196-3201};\\
Mohammad Arashi: \url{https://orcid.org/0000-0002-5881-9241}

\section{Introduction}
For the established two-periodic cyclical heavy-ball (CHB) coefficients of
Goujaud et al.~\cite[Algorithm 2 and Corollary 3.2]{goujaud2022}, we prove
an order-$R^4$ relative-error threshold as the normalized spectral gap $R$
closes at fixed condition number $\kappa>1$.

The exponent reflects a quadratic noiseless rate margin and square-root
sensitivity of repeated endpoint eigenvalues to scalar errors. The main
technical step is the matching dimension-independent lower bound for
history-dependent vector errors that mix eigenspaces. Unlike stationary
robustness results, we track this periodic family as its gap closes.
Switched-system certificates alone do not supply a matching obstruction;
we establish both bounds for the same coefficients, iterate-rate criterion
and ordinary first step from $x_{-1}=x_0$, without special rescaling. The comparison with stationary methods is a
corollary of the exact-gradient rate barrier; it does not identify the
best rate achievable after retuning for a positive error budget.

Table~\ref{tab:comparison} locates this contribution among robustness,
switched-system and energy analyses. Related schedules exploit spectral structure~\cite{oymak2021,agarwal2021}; broader-objective momentum is studied
in~\cite{khanh2026}. See also Vasin et al., \emph{On solving minimization
and min-max problems by first-order methods with relative error in gradients}
(2025, \href{https://arxiv.org/abs/2503.06628v3}{arXiv:2503.06628v3}).

\begin{table}[htbp]
\caption{Relation to the closest results. The comparisons distinguish the
coefficient choice and the rate criterion.}
\label{tab:comparison}
\small
\begin{tabular}{@{}>{\raggedright\arraybackslash}p{0.23\textwidth}>{\raggedright\arraybackslash}p{0.34\textwidth}>{\raggedright\arraybackslash}p{0.36\textwidth}@{}}
\toprule
Source & Established result & Contribution of this paper\\
\midrule
Goujaud et al.~\cite{goujaud2022} & Noiseless CHB rates on gapped spectra & Relative-error tolerance preserving the gap benefit\\[4pt]
Gannot~\cite{gannot2022}; Cyrus et al.~\cite{cyrus2018} & Relative-error rates and stationary rate--robustness tradeoffs & Fixed periodic coefficients; tolerance as the gap closes\\[4pt]
G\"urb\"uzbalaban~\cite{gurbuz2025} & Constant-momentum gains and multiplicative stability radii & Retention of $q$, with two-sided order in $R$\\[4pt]
Miller et al.~\cite[Theorem 4.1]{miller2026} & Switched convergence and rate certificates & Closed-form gap law and a matching scalar obstruction\\[4pt]
Orvieto~\cite[Proposition 4]{orvieto2025} & Invariant energy and convex-quadratic acceleration & Endpoint regularization for arbitrary vector errors\\[4pt]
De Klerk et al.~\cite[Theorem 5.1]{deklerk2017} & Tight noisy exact-line-search rate & Periodic fixed steps and an iterate-rate threshold\\
\botrule
\end{tabular}
\end{table}

\section{Problem, methods, and rates}\label{sec:setting}
Fix $0<\mu<L$, $\kappa=L/\mu>1$, and $0<R<1$.
For each positive integer $d$, consider $f:\R^d\to\R$ with
minimizer $x^\star\in\R^d$, minimum $f^\star\in\R$, and symmetric
Hessian $H\in\R^{d\times d}$ restricted to two spectral bands:
\begin{equation}\label{eq:class}
\begin{split}
 f(x)&=\tfrac12(x-x^\star)^TH(x-x^\star)+f^\star,\qquad
 \spec(H)\subseteq\Lambda_R=[\mu,a]\cup[b,L],\\
 a&=\frac{L+\mu-R(L-\mu)}2,\qquad
 b=\frac{L+\mu+R(L-\mu)}2 .
\end{split}
\end{equation}
In \eqref{eq:class}, $R=(b-a)/(L-\mu)$ is dimensionless and
$a-\mu=L-b$. The endpoints $\mu,L$ are class bounds, not necessarily attained
eigenvalues, and $a,b$ are the band edges defined in \eqref{eq:class}.
Scalar Hessians $H=\mu$ and $H=L$ are included.
We use the Euclidean and induced operator norms $\norm{\cdot}$,
transpose $T$, identity $I_n$ of order $n$, and
$A\preceq B$ for positive semidefinite $B-A$.

The CHB benchmark $q$, normalized midpoint $\chi$, momentum $m$,
step sizes $\alpha_0,\alpha_1$, and two-step margin $\Delta$ are
\begin{equation}\label{eq:coefficients}
\begin{gathered}
 q=\frac{\sqrt\kappa-1}{\sqrt\kappa+1},\qquad
 \chi=\frac{L+\mu}{L-\mu},\qquad
 m=\left(\frac{\sqrt{\chi^2-R^2}-\sqrt{\chi^2-1}}
                  {\sqrt{1-R^2}}\right)^2,\\
 \alpha_0=\frac{1+m}{a},\qquad
 \alpha_1=\frac{1+m}{b},\qquad
 \Delta=q^2-m .
\end{gathered}
\end{equation}
Using \eqref{eq:coefficients} from \cite[Algorithm 2]{goujaud2022}
and writing $\xi_k=x_k-x^\star$, we take one gradient query per step:
\begin{equation}\label{eq:chb}
\begin{split}
 x_{k+1}&=x_k+m(x_k-x_{k-1})-\alpha_r(H\xi_k+e_k),
 \qquad r=k\bmod 2,\\
 x_{-1}&=x_0,\qquad
 \norm{e_k}\leq\varepsilon\norm{H\xi_k},\qquad 0\leq\varepsilon<1 .
\end{split}
\end{equation}
For nonnegative integers $k$, the real errors $e_k\in\R^d$ may depend
adversarially on the entire history, without independence, unbiasedness,
symmetry, or eigenspace preservation. Equation \eqref{eq:chb} uses
$x_1=x_0-\alpha_0(H\xi_0+e_0)$ without rescaling and reports $x_k$.

For comparison, let $\mathcal S$ be the family with constant scalar
parameters $\alpha>0$ and $\beta,\eta\in\R$ defined by
\begin{equation}\label{eq:stationary}
\begin{split}
 y_k&=x_k+\eta(x_k-x_{k-1}),\\
 x_{k+1}&=x_k+\beta(x_k-x_{k-1})
       -\alpha\bigl(H(y_k-x^\star)+e_k\bigr),\\
 x_{-1}&=x_0,\qquad
 \norm{e_k}\leq\varepsilon\norm{H(y_k-x^\star)} .
\end{split}
\end{equation}
Parameters in \eqref{eq:stationary} may depend on $\mu,L,R,\varepsilon$,
but not individual eigenvalues. This includes gradient, heavy-ball,
Nesterov, and robust momentum, excluding matrix coefficients, higher
memory, and adaptation, with equal class information and one query per step.

For a fixed method $\mathcal A$ and class parameters,
$r_{\mathcal A}(\varepsilon)$ is the infimum of $r>0$ admitting
a finite constant $C$ such that
\begin{equation}\label{eq:rate}
 \norm{x_k-x^\star}\leq C r^k\norm{x_0-x^\star}
 \quad\text{for every }k\geq0,\ d,\ H,\ x_0,\text{ and admissible error path}.
\end{equation}
In \eqref{eq:rate}, $C$ may depend only on $\mu,L,R,\varepsilon$,
the method, and $r$. The empty-set infimum is $+\infty$;
the infimum need not be attained, notably at repeated roots.
Writing $r_{\mathcal C}$ for \eqref{eq:chb}, define the stationary
optimum and benchmark-rate threshold by
\begin{equation}\label{eq:threshold}
 r_{\mathcal S}^*(\varepsilon)
   =\inf_{\mathcal A\in\mathcal S}r_{\mathcal A}(\varepsilon),
 \qquad
 \varepsilon_{\crit}(\kappa,R)
   =\sup\bigl(\{0\}\cup
     \{\varepsilon\in[0,1):r_{\mathcal C}(\varepsilon)<q\}\bigr).
\end{equation}
Scaling $H,\mu,L$ leaves $\alpha_rH$ unchanged, so
\eqref{eq:threshold} depends only on $\kappa,R$. Nested error sets make
$r_{\mathcal C}(\varepsilon)$ nondecreasing. The threshold retains $q$;
it measures neither divergence nor the full rate--robustness tradeoff
of methods retuned for inexact gradients.

\section{An explicit fourth-power law}\label{sec:main}
The following constants depend only on $\kappa$. The quantities $d_-$
and $d_+$ bound $\Delta/R^2$. For $R\leq1/2$ and phase $r\in\{0,1\}$,
$\mathsf A_r$ bounds $\alpha_rL$, $\mathsf B_r$ bounds
$|1+m-\alpha_r\lambda|$, and $\cT_r$ bounds the state transition;
$\cK$ and $\cG$ control the energy and two-step perturbation:
\begin{equation}\label{eq:auxconstants}
\begin{gathered}
 d_0=\frac{2q^2\sqrt\kappa}{\kappa+1},\qquad
 d_-=\tfrac34d_0,\qquad d_+=2d_0,\\
 \mathsf A_0=\frac{8\kappa}{\kappa+3},\quad
 \mathsf A_1=\frac{4\kappa}{\kappa+1},\quad
 \mathsf B_0=\frac{6(\kappa-1)}{\kappa+3},\quad
 \mathsf B_1=\frac{6(\kappa-1)}{3\kappa+1},\\
 \cT_r=\sqrt{\mathsf B_r^2+2},\qquad
 \cK=\sqrt{(\cT_0\cT_1+1)^2+2},\qquad
 \cG=\cT_1\mathsf A_0+\cT_0\mathsf A_1+\mathsf A_0\mathsf A_1.
\end{gathered}
\end{equation}
Using \eqref{eq:auxconstants}, choose the positive threshold constants
and the admissible gap range as follows:
\begin{equation}\label{eq:mainconstants}
\begin{gathered}
 c_\kappa=\min\left\{1,\frac{d_-^2}{4\cK\cG}\right\},\qquad
 C_\kappa=\frac{4q^2}{(\kappa-1)(1+3q^2/4)^2},\\
 R_\kappa=\min\left\{\frac12,(2C_\kappa)^{-1/4}\right\}.
\end{gathered}
\end{equation}

\begin{theorem}\label{thm:main}
For every fixed $\kappa>1$ and $0<R\leq R_\kappa$, the method
\eqref{eq:chb} and the threshold \eqref{eq:threshold} satisfy
\begin{equation}\label{eq:fourthpower}
 c_\kappa R^4\leq\varepsilon_{\crit}(\kappa,R)\leq C_\kappa R^4 .
\end{equation}
For $0\leq\varepsilon<c_\kappa R^4$, set $\delta=\Delta/2$.
An explicit bound for every iterate, implying
$r_{\mathcal C}(\varepsilon)<q$, is
\begin{equation}\label{eq:prefactor}
\begin{split}
 \gamma&=m+\delta+\frac{\cK\cG}{\delta}\varepsilon<q^2,\\
 \mathcal C&=\frac{\sqrt2\cK}{\delta}
     \max\left\{1,\frac{\cT_0+\mathsf A_0\varepsilon}{\sqrt\gamma}\right\},\\
 \norm{\xi_k}&\leq\mathcal C\,\gamma^{k/2}\norm{\xi_0}
 \quad(k\geq0).
\end{split}
\end{equation}
The constants in \eqref{eq:prefactor} are finite, positive, and
independent of dimension, Hessian, initial point, and error path.
\end{theorem}

The bounds in \eqref{eq:fourthpower} have matching fourth-power order
at fixed $\kappa$; their coefficients are not claimed optimal.
No uniform assertion is made as $\kappa\to\infty$ and $R\to0$
simultaneously. The prefactor in \eqref{eq:prefactor} can grow as the
gap closes, so the rate bound alone implies no short-horizon advantage.

\begin{corollary}\label{cor:comparison}
Under the assumptions of Theorem~\ref{thm:main}, if
$0\leq\varepsilon<c_\kappa R^4$, then
\begin{equation}\label{eq:separation}
 r_{\mathcal C}(\varepsilon)<q=r_{\mathcal S}^*(0)
     \leq r_{\mathcal S}^*(\varepsilon).
\end{equation}
\end{corollary}
\begin{proof}
Theorem~\ref{thm:main} gives the strict inequality in
\eqref{eq:separation}; Lemma~\ref{lem:stationary} below gives the rest.
\end{proof}
This comparison concerns only $\mathcal S$. It uses the admissibility
of zero error, not a strict deterioration of the optimal stationary
rate when the error budget increases.

\section{Proof of Theorem \ref{thm:main}}\label{sec:proof}
\subsection{The stationary benchmark and initialization}
\begin{lemma}\label{lem:stationary}
For every $\varepsilon\in[0,1)$, the family
\eqref{eq:stationary} satisfies
$r_{\mathcal S}^*(\varepsilon)\geq r_{\mathcal S}^*(0)=q$.
\end{lemma}
\begin{proof}
Zero error is admissible for every $\varepsilon$. Take a scalar Hessian
$H=\lambda$, where $\lambda$ is either $\mu$ or $L$, and set
$x^\star=0$, $x_{-1}=x_0=1$. The state matrix and its characteristic
polynomial, with scalar variable $z$, are
\begin{equation}\label{eq:stationarymatrix}
\begin{split}
 S_\lambda&=
 \begin{pmatrix}
 1+\beta-\alpha\lambda(1+\eta)&-\beta+\alpha\eta\lambda\\
 1&0
 \end{pmatrix},\\
 p_\lambda(z)&=z^2-[1+\beta-\alpha\lambda(1+\eta)]z
                  +\beta-\alpha\eta\lambda,\qquad
 p_\lambda(1)=\alpha\lambda .
\end{split}
\end{equation}
For the initialized state $v=(1,1)^T$, \eqref{eq:stationarymatrix}
gives $S_\lambda v=(1-\alpha\lambda,1)^T$ and
$\det[v,S_\lambda v]=\alpha\lambda>0$. Thus $v,S_\lambda v$
form a basis, regardless of the parameters $\beta,\eta$.

Suppose \eqref{eq:rate} held with some $0<r<q$.
The bounds on $x_k,x_{k-1}$ imply
$\norm{S_\lambda^kv}\leq C_\lambda r^k$, enlarging the constant to
include $k=0$. They also bound $S_\lambda^k(S_\lambda v)$ by a
constant times $r^k$. Using the basis just established gives
$\norm{S_\lambda^k}\leq C'_\lambda r^k$. Consequently both
characteristic roots $z_1,z_2$ have modulus at most $r$.
This argument rules out cancellation of a root by the prescribed
initialization.

Evaluating the polynomial in \eqref{eq:stationarymatrix} at $1$
now bounds the step size at both endpoints:
\begin{equation}\label{eq:endpointobstruction}
 (1-r)^2\leq |(1-z_1)(1-z_2)|=\alpha\lambda\leq(1+r)^2,
 \qquad
 \kappa=\frac{\alpha L}{\alpha\mu}
       \leq\left(\frac{1+r}{1-r}\right)^2.
\end{equation}
Since $r<1$, \eqref{eq:endpointobstruction} implies $r\geq q$,
a contradiction. If $r_{\mathcal A}(\varepsilon)<q$, the infimum
definition would supply an admissible $r<q$. Therefore no member of
$\mathcal S$ has a rate below $q$.

For the reverse bound at zero error, stationary heavy ball with
$\alpha=4/(\sqrt L+\sqrt\mu)^2$, $\beta=q^2$ and $\eta=0$
has $\det S_\lambda=q^2$ and $|\tr S_\lambda|\leq2q$
for every $\lambda\in[\mu,L]$. Lemma~\ref{lem:energy} below,
with $m=q$ and any $\delta>0$, bounds its scalar powers by
$C_\delta(q+\delta)^k$ uniformly over this compact interval:
the energy has lower bound $\delta^2I_2$ and a uniform upper
bound because the entries of $S_\lambda$ are bounded.
Taking direct sums and using $\norm{z_0}=\sqrt2\norm{\xi_0}$
preserves uniformity in dimension and initialization. Letting
$\delta\downarrow0$ in the rate infimum gives
$r_{\mathcal S}^*(0)\leq q$. Finally, nested error sets give
$r_{\mathcal A}(\varepsilon)\geq r_{\mathcal A}(0)$ for every
fixed parameter triple. Taking infima over the same triples proves
$r_{\mathcal S}^*(\varepsilon)\geq r_{\mathcal S}^*(0)$, even
when the optimizing triple depends on $\varepsilon$.
\end{proof}

\subsection{The size of the noiseless margin}
\begin{lemma}\label{lem:margin}
For $0<R\leq1/2$, the coefficients \eqref{eq:coefficients} satisfy
$0<m<q^2<1$ and $d_-R^2\leq\Delta\leq d_+R^2$.
\end{lemma}
\begin{proof}
Let $u=R^2$, $B=\sqrt{\chi^2-1}>0$, and
$A(u)=\sqrt{\chi^2-u}$. Rationalizing the numerator in
\eqref{eq:coefficients} yields the following representation of
$s(u)=\sqrt{m(u)}$:
\begin{equation}\label{eq:rationalm}
 s(u)=\frac{\sqrt{1-u}}{A(u)+B},\qquad
 s(0)=\frac1{\chi+B}=q,\qquad
 m(u)=\frac{1-u}{(A(u)+B)^2}.
\end{equation}
Taking logarithmic derivatives in \eqref{eq:rationalm} and using
$A(u)^2-B^2=1-u$ gives exact derivative identities:
\begin{equation}\label{eq:derivative}
 \frac{s'(u)}{s(u)}=-\frac{B}{2A(u)(1-u)},\qquad
 -m'(u)=\frac{m(u)B}{A(u)(1-u)},\qquad
 -m'(0)=\frac{q^2B}{\chi}=d_0.
\end{equation}
The last equality in \eqref{eq:derivative} follows from
$B/\chi=2\sqrt\kappa/(\kappa+1)$.

For $0\leq u\leq1/4$, \eqref{eq:rationalm} implies
$m(u)\geq(1-u)q^2\geq3q^2/4$, while
\eqref{eq:derivative} implies $m(u)\leq q^2$.
Also $3\chi/4\leq A(u)\leq\chi$: for the lower bound,
$\chi^2-u\geq\chi^2-1/4\geq9\chi^2/16$ because $\chi>1$.
Together with $1-u\geq3/4$, these estimates give
$(3/4)d_0\leq -m'(u)\leq(16/9)d_0<2d_0$.
Integration therefore provides the required explicit margin:
\begin{equation}\label{eq:margin}
 \Delta=\int_0^{R^2}[-m'(u)]\,du,\qquad
 d_-R^2\leq\Delta\leq d_+R^2 .
\end{equation}
Equation \eqref{eq:margin} also proves $\Delta>0$. These bounds hold
throughout the stated interval.
\end{proof}

\subsection{The two-step transition}
\begin{lemma}\label{lem:trace}
For every $\lambda\in\Lambda_R$, the noiseless two-step state matrix
has determinant $m^2$ and trace of absolute value at most $2m$.
\end{lemma}
\begin{proof}
For phase $r\in\{0,1\}$ put $c_r(\lambda)=1+m-\alpha_r\lambda$.
The one-step and two-step matrices are
\begin{equation}\label{eq:transitions}
 T_r(\lambda)=\begin{pmatrix}c_r(\lambda)&-m\\1&0\end{pmatrix},
 \qquad
 M_\lambda=T_1(\lambda)T_0(\lambda)
   =\begin{pmatrix}c_1c_0-m&-mc_1\\c_0&-m\end{pmatrix}.
\end{equation}
Each factor in \eqref{eq:transitions} has determinant $m$, so
$\det M_\lambda=m^2$. Substitution of the step sizes gives
the trace polynomial
\begin{equation}\label{eq:tracepoly}
 \tr M_\lambda
   =(1+m)^2\frac{(\lambda-a)(\lambda-b)}{ab}-2m .
\end{equation}
To evaluate \eqref{eq:tracepoly} at the band endpoints, use the
quantities $A,B,u,s$ from Lemma \ref{lem:margin}.
The identities
$s=(A-B)/\sqrt{1-u}$ and $1/s=(A+B)/\sqrt{1-u}$ imply
$s+1/s=2A/\sqrt{1-u}$. Squaring this equality gives
\begin{equation}\label{eq:equioscillation}
 (1+m)^2\frac{1-R^2}{\chi^2-R^2}=4m,\qquad
 \frac{(L-a)(L-b)}{ab}
 =\frac{(\mu-a)(\mu-b)}{ab}
 =\frac{1-R^2}{\chi^2-R^2}.
\end{equation}
Equations \eqref{eq:tracepoly} and \eqref{eq:equioscillation}
give trace $2m$ at $\mu,L$ and trace $-2m$ at $a,b$.
The polynomial $(\lambda-a)(\lambda-b)$ decreases on $[\mu,a]$
and increases on $[b,L]$, since its derivative is $2\lambda-a-b$
and $a+b=L+\mu$. Thus its values do not exceed the endpoint
range within either band. This proves the trace bound.
\end{proof}

\subsection{An energy valid at repeated roots}
For a real symmetric positive definite matrix $P$, use the energy
norm $\norm{v}_P=(v^TPv)^{1/2}$ and its induced operator norm
$\norm{A}_P=\sup_{v\ne0}\norm{Av}_P/\norm{v}_P$.
The next lemma avoids an eigenvector condition number, which would
be singular at a defective endpoint.

\begin{lemma}\label{lem:energy}
Let $M\in\R^{2\times2}$, $m>0$, $\det M=m^2$, and
$|\tr M|\leq2m$. For every $\delta>0$, there is a symmetric
$P_\delta\succeq\delta^2I_2$ such that
$\norm{M}_{P_\delta}\leq m+\delta$.
\end{lemma}
\begin{proof}
Set $t=\tr M/2$, $N=M-tI_2$, and $\omega^2=m^2-t^2\geq0$.
Cayley--Hamilton gives $N^2=-\omega^2I_2$. Define the
unregularized and regularized energies by
\begin{equation}\label{eq:energydef}
 P_0=N^TN+\omega^2I_2,\qquad
 P_\delta=P_0+\delta^2I_2 .
\end{equation}
The first matrix in \eqref{eq:energydef} can be singular.
The second is positive definite even when $\omega=0$.
Direct multiplication using $N^2=-\omega^2I_2$ yields
\begin{equation}\label{eq:energyidentities}
 N^TP_0+P_0N=0,\qquad
 N^TP_0N=\omega^2P_0,\qquad
 M^TP_0M=(t^2+\omega^2)P_0=m^2P_0 .
\end{equation}
For example, in the first identity in \eqref{eq:energyidentities},
$(N^T)^2N=-\omega^2N$ and $N^TN^2=-\omega^2N^T$ cancel the
remaining terms. In the second, $(N^2)^TN^2=\omega^4I_2$.
Expanding the regularization using these identities gives an
explicit sum of positive semidefinite terms:
\begin{equation}\label{eq:sos}
\begin{split}
 &(m^2+m\delta+\delta^2)P_\delta-M^TP_\delta M\\
 &\quad=m\delta
       \left(N-\frac{t\delta}{m}I_2\right)^T
       \left(N-\frac{t\delta}{m}I_2\right)\\
 &\qquad\quad+
   \left[\omega^2\left(m\delta+2\delta^2+\frac{\delta^3}{m}\right)
                 +\delta^4\right]I_2 .
\end{split}
\end{equation}
Every coefficient on the right of \eqref{eq:sos} is nonnegative.
Consequently the desired matrix inequality follows:
\begin{equation}\label{eq:energybound}
 M^TP_\delta M\preceq(m^2+m\delta+\delta^2)P_\delta
                  \preceq(m+\delta)^2P_\delta .
\end{equation}
By the definition of the induced norm, \eqref{eq:energybound}
proves the claim. Neither diagonalizability nor a nonzero
$\omega$ was assumed.
\end{proof}

\subsection{Uniform control of arbitrary vector errors}
\begin{lemma}\label{lem:vector}
Let $0<R\leq1/2$, $0<\delta\leq1$, and $0\leq\varepsilon<1$.
For each admissible Hessian $H$, a single positive definite energy
matrix $P_H$ bounds every two-step block of \eqref{eq:chb}
by $m+\delta+\cK\cG\varepsilon/\delta$ in its induced norm.
The norm-equivalence constants are independent of $H$ and $d$.
\end{lemma}
\begin{proof}
Let $H=Q\diag(\lambda_1,\ldots,\lambda_d)Q^T$, where $Q$ is
orthogonal. In the state $z_k=(\xi_k^T,\xi_{k-1}^T)^T$, apply
$Q^T$ to each component and permute the coordinates into pairs.
The noiseless two-step matrix $M_H=T_1(H)T_0(H)$ then becomes
the direct sum of the matrices $M_{\lambda_i}$ in
\eqref{eq:transitions}. Here
$T_r(H)=\left(\begin{smallmatrix}(1+m)I_d-\alpha_rH&-mI_d\\
I_d&0\end{smallmatrix}\right)$.
Take the direct sum of the energies in Lemma \ref{lem:energy}
and transform it back by the same orthogonal change of coordinates;
call the resulting matrix $P_H$.

We bound the phases separately. From $0<R\leq1/2$ and
\eqref{eq:class},
$(L+3\mu)/4\leq a\leq(L+\mu)/2$ and
$(L+\mu)/2\leq b\leq(3L+\mu)/4$.
Together with $1+m<2$, these intervals imply
$\alpha_rL\leq\mathsf A_r$. For $\lambda\in[\mu,L]$,
$1-\lambda/a$ lies between $-3(\kappa-1)/(\kappa+3)$ and
$(\kappa-1)/(\kappa+1)$; $1-\lambda/b$ lies between
$-(\kappa-1)/(\kappa+1)$ and $3(\kappa-1)/(3\kappa+1)$.
The larger absolute endpoint in each interval gives
$|c_r(\lambda)|\leq\mathsf B_r$.
The scalar Frobenius norm therefore satisfies
$\norm{T_r(\lambda)}\leq\sqrt{\mathsf B_r^2+m^2+1}\leq\cT_r$.
The norm of the full direct sum is its largest block norm, so
$\norm{T_r(H)}\leq\cT_r$ without a dimension factor.

In each scalar block, $|t|\leq m\leq1$,
$\norm{N}\leq\norm{M_\lambda}+|t|\leq\cT_1\cT_0+1$, and
$\omega^2\leq1$. Since $\delta\leq1$,
\eqref{eq:energydef} and Lemma \ref{lem:energy} yield
\begin{equation}\label{eq:uniformenergy}
\begin{gathered}
 \delta^2I_{2d}\preceq P_H\preceq\cK^2I_{2d},\qquad
 \norm{M_H}_{P_H}\leq m+\delta,\\
 \norm{F}_{P_H}\leq\frac{\cK}{\delta}\norm{F}
 \qquad\text{for every }F\in\R^{2d\times2d}.
\end{gathered}
\end{equation}
The choice $\cK^2=(\cT_0\cT_1+1)^2+2$ in
\eqref{eq:auxconstants} proves the upper bound in
\eqref{eq:uniformenergy}. The last inequality follows from its
two norm-equivalence bounds. It applies to all matrices $F$,
including those that mix the Hessian eigenspaces. Although $P_H$
depends on $H$, the constants in \eqref{eq:uniformenergy} do not.

To represent an arbitrary error along a realized path, let
$h_k=H\xi_k$. If $h_k\ne0$, define a real rank-one matrix by
\begin{equation}\label{eq:errorrepresentation}
 D_k=\frac{e_kh_k^T}{\norm{h_k}^2},\qquad
 e_k=D_kh_k,\qquad
 \norm{D_k}=\frac{\norm{e_k}}{\norm{h_k}}\leq\varepsilon .
\end{equation}
If $h_k=0$, the relative error assumption forces $e_k=0$ and
we set $D_k=0$. Equation \eqref{eq:errorrepresentation} is
a pathwise representation, not an assumption on the error mechanism.
In particular $D_k$ may depend on the current state and the history.
For $k=2j+r$, the perturbed state transition is
\begin{equation}\label{eq:onesteppert}
\begin{split}
 z_{k+1}&=(T_r(H)+E_{r,j})z_k,\\
 E_{r,j}&=\begin{pmatrix}-\alpha_rD_{2j+r}H&0\\0&0\end{pmatrix},
 \qquad \norm{E_{r,j}}\leq\mathsf A_r\varepsilon .
\end{split}
\end{equation}
All blocks in \eqref{eq:onesteppert} have size $d\times d$.
Multiplying two consecutive transitions along the same path gives
\begin{equation}\label{eq:twosteppert}
\begin{split}
 z_{2j+2}&=(M_H+F_j)z_{2j},\\
 F_j&=T_1(H)E_{0,j}+E_{1,j}T_0(H)+E_{1,j}E_{0,j},\\
 \norm{F_j}&\leq(\cT_1\mathsf A_0+\cT_0\mathsf A_1)\varepsilon
                  +\mathsf A_0\mathsf A_1\varepsilon^2
             \leq\cG\varepsilon .
\end{split}
\end{equation}
The final bound in \eqref{eq:twosteppert} uses $\varepsilon<1$
and the definition of $\cG$. Dependence of $D_{2j+1}$ on the first
step does not alter this exact pathwise product.

The triangle inequality, together with \eqref{eq:uniformenergy} and
\eqref{eq:twosteppert}, bounds every block in the same energy:
\begin{equation}\label{eq:blockbound}
 \norm{M_H+F_j}_{P_H}\leq
 \gamma:=m+\delta+\frac{\cK\cG}{\delta}\varepsilon,\qquad
 \norm{z_{2j}}\leq\frac{\cK}{\delta}\gamma^j\norm{z_0}.
\end{equation}
The second estimate in \eqref{eq:blockbound} follows by iterating in
the energy norm and converting to the Euclidean norm only at the initial
and final states. Thus the norm-equivalence constants enter once.

Initially, $\norm{z_0}=\sqrt2\norm{\xi_0}$.
For odd indices we use the bound
$\norm{T_0(H)+E_{0,j}}\leq\cT_0+\mathsf A_0\varepsilon$.
Equation \eqref{eq:blockbound} then gives both parity bounds:
\begin{equation}\label{eq:parities}
\begin{split}
 \norm{\xi_{2j}}&\leq
       \frac{\sqrt2\cK}{\delta}\gamma^j\norm{\xi_0},\\
 \norm{\xi_{2j+1}}&\leq
       \frac{\sqrt2\cK}{\delta}(\cT_0+\mathsf A_0\varepsilon)
                    \gamma^j\norm{\xi_0}.
\end{split}
\end{equation}
Since $\gamma>0$, the two estimates in \eqref{eq:parities}
combine into the prefactor form of \eqref{eq:prefactor}.
This proves the lemma, including its uniformity assertion.
\end{proof}

\subsection{The sufficient error level}
We prove the lower bound in \eqref{eq:fourthpower} and the
all-iterate estimate \eqref{eq:prefactor}.
Lemma \ref{lem:margin} gives $0<\Delta<q^2<1$, so
$\delta=\Delta/2$ is permitted in Lemma \ref{lem:vector}.
For $0\leq\varepsilon<c_\kappa R^4$, we have
$\varepsilon<1$ because $c_\kappa\leq1$ and $R\leq1/2$.
The choices in \eqref{eq:mainconstants} and the margin
\eqref{eq:margin} imply
\begin{equation}\label{eq:strictmargin}
 \frac{\cK\cG\varepsilon}{\delta}
 <\frac{d_-^2R^4}{4\delta}
 \leq\frac{\Delta^2}{4\delta}
 =\frac{\Delta}{2},\qquad
 \gamma<m+\frac{\Delta}{2}+\frac{\Delta}{2}=q^2 .
\end{equation}
Using \eqref{eq:strictmargin} in \eqref{eq:parities} proves
\eqref{eq:prefactor} and
$r_{\mathcal C}(\varepsilon)\leq\sqrt\gamma<q$.
Every $\varepsilon<c_\kappa R^4$ is in the strict-rate set
of \eqref{eq:threshold}, so its supremum is at least
$c_\kappa R^4$. No strict-rate claim is needed at
the endpoint $\varepsilon=c_\kappa R^4$.

\subsection{A matching upper order and output observability}
\begin{lemma}\label{lem:upper}
Let $0<R\leq1/2$. Define positive dimensionless quantities
$D_R,B_R,V$ and a scalar upper threshold $\varepsilon_U$ by
\begin{equation}\label{eq:upperthreshold}
\begin{gathered}
 D_R=\alpha_0\alpha_1L(L-\mu),\qquad
 B_R=\alpha_0\alpha_1L^2,\qquad
 V=\frac{\Delta^2}{q^2},\\
 \varepsilon_U
 =\frac{\sqrt{D_R^2+4B_RV}-D_R}{2B_R}
 =\frac{2V}{D_R+\sqrt{D_R^2+4B_RV}} .
\end{gathered}
\end{equation}
If $\varepsilon_U<1$, then
$\varepsilon_{\crit}(\kappa,R)\leq\varepsilon_U$.
\end{lemma}
\begin{proof}
Take $d=1$, $H=L$, $x^\star=0$, $x_{-1}=x_0=1$,
and the admissible multiplicative error $e_k=\varepsilon Lx_k$.
The actual Hessian remains $L$ and satisfies \eqref{eq:class}.
The state transition is obtained from \eqref{eq:transitions}
at the effective value $\lambda_{\mathrm{eff}}=L(1+\varepsilon)$.
This substitution describes the erroneous oracle, not a Hessian
outside the problem class.

The trace formula \eqref{eq:tracepoly}, $a+b=L+\mu$, and
$\tr M_L=2m$ give the perturbed trace and its larger eigenvalue:
\begin{equation}\label{eq:upperroot}
\begin{split}
 \tr M_{\lambda_{\mathrm{eff}}}&=2m+u_\varepsilon,\qquad
 u_\varepsilon=D_R\varepsilon+B_R\varepsilon^2,\\
 \rho_+(\varepsilon)&=
       m+\frac{u_\varepsilon}{2}
        +\sqrt{mu_\varepsilon+\frac{u_\varepsilon^2}{4}} .
\end{split}
\end{equation}
The determinant remains $m^2$. For $\varepsilon>0$ the two
eigenvalues are distinct and positive. The larger one in
\eqref{eq:upperroot} equals $q^2$ precisely when the trace is
$q^2+m^2/q^2$, or equivalently when
$u_\varepsilon=(q^2-m)^2/q^2=V$.
The unique positive solution is \eqref{eq:upperthreshold}.
Both $u_\varepsilon$ and $\rho_+(\varepsilon)$ increase with
$\varepsilon$, so $\rho_+(\varepsilon)\geq q^2$ for
$\varepsilon\geq\varepsilon_U>0$.

It remains to check that this larger eigenvalue appears in the
initialized output. At such an $\varepsilon$,
$\lambda_{\mathrm{eff}}>L>b>a$, hence both
$c_0(\lambda_{\mathrm{eff}})$ and
$c_1(\lambda_{\mathrm{eff}})$ are negative.
For $v=(1,1)^T$ and $M=M_{\lambda_{\mathrm{eff}}}$,
$(Mv)_2=c_0-m<0$. If $v$ were an eigenvector for the smaller
positive eigenvalue, this component would be positive.
Thus the coefficient of the larger eigenvector in the expansion
of $v$ is nonzero. Further, $M_{12}=-mc_1>0$, so no eigenvector
can have zero first coordinate. The larger eigenvalue is therefore
both excited by $v$ and visible in $x_{2j}$. Its contribution gives
the exact output rate
\begin{equation}\label{eq:observability}
 \lim_{j\to\infty}|x_{2j}|^{1/(2j)}
     =\sqrt{\rho_+(\varepsilon)}\geq q
 \qquad(\varepsilon_U\leq\varepsilon<1).
\end{equation}
Equation \eqref{eq:observability} excludes every such
$\varepsilon$ from the strict-rate set in \eqref{eq:threshold}.
This proves the lemma with the prescribed initialization and
the actual reported output.
\end{proof}

To finish the upper bound in Theorem \ref{thm:main}, note that
$D_R\varepsilon_U+B_R\varepsilon_U^2=V$ implies
$\varepsilon_U\leq V/D_R$.
Since $ab\leq(L+\mu)^2/4$ and $m\geq3q^2/4$ by
Lemma~\ref{lem:margin}, we have the
explicit lower bound on the linear trace perturbation:
\begin{equation}\label{eq:Dlower}
 D_R=\frac{(1+m)^2L(L-\mu)}{ab}
 \geq\frac{4(1+3q^2/4)^2L(L-\mu)}{(L+\mu)^2}
 =\frac{4(1+3q^2/4)^2\kappa(\kappa-1)}{(\kappa+1)^2}.
\end{equation}
Combining \eqref{eq:Dlower} with Lemma \ref{lem:margin} gives
the required fourth-power upper bound:
\begin{equation}\label{eq:upperconstant}
 \varepsilon_U\leq\frac{\Delta^2}{q^2D_R}
 \leq\frac{d_+^2}{q^2}
         \frac{(\kappa+1)^2}{4(1+3q^2/4)^2\kappa(\kappa-1)}R^4
 =C_\kappa R^4 .
\end{equation}
For $R\leq R_\kappa$, \eqref{eq:mainconstants} ensures
$C_\kappa R^4\leq1/2$. Thus \eqref{eq:upperconstant} puts
$\varepsilon_U$ strictly inside the permitted error interval,
and Lemma \ref{lem:upper} applies.
This proves the upper bound in \eqref{eq:fourthpower}.
The lower bound and the all-iterate estimate were proved above,
completing Theorem \ref{thm:main}. \hfill$\square$
\section{Two analytic refinements}\label{sec:refinements}
The first refinement identifies the noiseless rate for our exact
initialization. It also quantifies the gain that the error threshold
preserves.

\begin{corollary}\label{cor:noiseless}
For $0<R\leq1/2$, $r_{\mathcal C}(0)=\sqrt m$.
The noiseless rate advantage obeys
\begin{equation}\label{eq:noiselessgain}
 \frac{d_-}{2q}R^2\leq q-\sqrt m\leq\frac{d_+}{q}R^2,
 \qquad
 \lim_{R\downarrow0}\frac{q-\sqrt m}{R^2}
       =\frac{q\sqrt\kappa}{\kappa+1}.
\end{equation}
\end{corollary}
\begin{proof}
At zero error, Lemma \ref{lem:vector} applies for every sufficiently
small $\delta>0$, giving
$r_{\mathcal C}(0)\leq\sqrt{m+\delta}$.
Letting $\delta\downarrow0$ proves
$r_{\mathcal C}(0)\leq\sqrt m$.
The prefactor need not remain bounded in this limit: the rate is
an infimum as defined in \eqref{eq:rate}.

For the reverse inequality take $H=L$, $x^\star=0$, and
$x_{-1}=x_0=1$. Lemma \ref{lem:trace} gives trace $2m$ and
determinant $m^2$. With $N=M_L-mI_2$,
Cayley--Hamilton yields $N^2=0$, and therefore the exact power
formula is
\begin{equation}\label{eq:jordanpower}
 M_L^j=m^j\left(I_2+\frac{j}{m}N\right),\qquad j\geq0 .
\end{equation}
For any nonzero vector $v$, the factor in parentheses in
\eqref{eq:jordanpower} sends $v$ either to $v$ (if $Nv=0$)
or to a vector with norm asymptotically linear in $j$.
Thus $\norm{M_L^jv}^{1/j}\to m$.
An output bound \eqref{eq:rate} with $r<\sqrt m$ would bound
both coordinates of $z_{2j}=(x_{2j},x_{2j-1})^T$ by a
constant times $r^{2j}$, contradicting this limit.
This establishes the claimed noiseless rate.

Finally,
$\Delta=(q-\sqrt m)(q+\sqrt m)$ and
$q\leq q+\sqrt m\leq2q$ convert Lemma \ref{lem:margin}
to the bounds in \eqref{eq:noiselessgain}.
Equation \eqref{eq:derivative} gives $\Delta/R^2\to d_0$;
dividing by $q+\sqrt m\to2q$ proves the limit in
\eqref{eq:noiselessgain}.
\end{proof}

The second refinement concerns the explicit scalar upper threshold,
not the exact worst-case threshold over all vector errors.

\begin{corollary}\label{cor:uppercoefficient}
For fixed $\kappa>1$, the upper threshold
\eqref{eq:upperthreshold} satisfies
\begin{equation}\label{eq:upperlimit}
 \lim_{R\downarrow0}\frac{\varepsilon_U}{R^4}
      =\frac{\kappa-1}{4(\kappa+1)^2}.
\end{equation}
\end{corollary}
\begin{proof}
As $R\downarrow0$, the coefficient formula
\eqref{eq:coefficients} gives
$m\to q^2$ and
$\alpha_0,\alpha_1\to4/[\mu(\sqrt\kappa+1)^2]$.
Consequently
$D_R\to D_0=16\kappa(\kappa-1)/(\sqrt\kappa+1)^4>0$,
while $B_R$ remains bounded. By \eqref{eq:derivative},
$\Delta/R^2\to d_0$.
The rationalized expression in \eqref{eq:upperthreshold}
therefore gives
$\varepsilon_U/R^4\to d_0^2/(q^2D_0)$.
Substituting the values of $d_0,q,D_0$ simplifies this to
\eqref{eq:upperlimit}.
\end{proof}

The scalar construction is only one member of the allowed error
class. Thus \eqref{eq:upperlimit} does not assert the same limiting
coefficient for $\varepsilon_{\crit}$. Errors that vary over time or
mix eigenspaces might produce a smaller critical constant.

\section{Controlled two-band examples}\label{sec:numerics}
\suppressfloats[t]
Set $\mu=1/16$, $L=1$, $\kappa=16$, and $q=0.6$ for scalar
thresholds and 512-query comparisons with fixed noiseless tuning.
These controlled quadratics are the primary numerical examples: they
isolate the two-band geometry at a moderate condition number.
The case $R=1/2$ illustrates the spectral-gap benefit, while $R=1/8$
shows its reduction as the gap narrows. The two cases keep the
condition number, initialization, and query budget fixed.

\subsection{Scalar threshold scaling}
Figure~\ref{fig:scaling} evaluates \eqref{eq:upperthreshold} at
$R=2^{-j}$, $j=1,\ldots,5$. The local exponent
$\log(\varepsilon_U(2R)/\varepsilon_U(R))/\log 2$, plotted at the
smaller gap, decreases from $4.114858$ to $4.001647$.
At $R=1/32$, $\varepsilon_U/R^4=0.0129807135$ is $0.0380\%$
above the limit $15/1156$ in \eqref{eq:upperlimit}.
For $H=L$, the per-query factors at $\varepsilon_U/2$, $\varepsilon_U$,
and $3\varepsilon_U/2$ lie below, at, and above $q$, respectively.
This illustrates the scalar obstruction, not the vector-error threshold.

\begin{figure}[!htbp]
\centering
\includegraphics[width=\textwidth]{figures_v24/Fig1.pdf}
\caption{Scalar upper obstruction at $\kappa=16$. The left panel
normalizes \eqref{eq:upperthreshold} by $R^4$; the horizontal line
is the limiting coefficient in \eqref{eq:upperlimit}. The right
panel shows the local exponent for consecutive gap halvings.}
\label{fig:scaling}
\end{figure}

\subsection{Method comparison at a fixed query budget}
For $R\in\{1/2,1/8\}$, use $H=\diag(\mu,a,b,L)$ from
\eqref{eq:class}, $x^\star=0$, and $x_{-1}=x_0=(1,1,1,1)^T$.
In the stationary family \eqref{eq:stationary}, HB uses $\eta=0$,
$\alpha=4/(\sqrt L+\sqrt\mu)^2$, and $\beta=q^2$.
Gradient descent (GD) uses $\alpha=2/(L+\mu)$ and $\beta=\eta=0$.
The quadratically tuned Nesterov method, denoted NAG-Q, uses
$\alpha=4/(3L+\mu)$ and
$\beta=\eta=(\sqrt{3\kappa+1}-2)/(\sqrt{3\kappa+1}+2)$
\cite[Proposition~1]{lessard2016}. Its gradient, including the
relative error, is evaluated at $y_k$ in \eqref{eq:stationary}.
All four methods share Hessians, initializations, error policies,
and query counts; none is retuned under errors.

The four-method comparison uses exact gradients and scalar errors
$e_k=(\varepsilon_U/2)\nabla f(y_k)$, with $y_k=x_k$ for CHB, HB,
and GD. Figure~\ref{fig:finite} plots $D_k=\norm{x_k}/\norm{x_0}$.
CHB has the smallest terminal distance in both problems and policies,
also when comparing $\max\{D_{511},D_{512}\}$. With exact gradients,
the terminal CHB/HB distance ratio increases from
$9.9911\times10^{-17}$ at $R=1/2$ to $0.1383215$ at $R=1/8$:
the advantage over HB weakens as the gap narrows.
The corresponding CHB/NAG-Q ratios are $9.7639\times10^{-55}$
and $1.3518\times10^{-39}$.
\begin{figure}[!htbp]
\centering
\includegraphics[width=\textwidth]{figures_v24/Fig2.pdf}
\caption{Normalized iterate distances on the controlled two-band quadratics,
with rest initialization and all four endpoint modes excited.
The left column is the primary $R=1/2$ example; the right column
reduces the gap to $R=1/8$ at the same condition number $\kappa=16$.
The top row uses exact gradients; the bottom row uses scalar
relative errors with $\varepsilon=\varepsilon_U/2$.
Each method retains its fixed noiseless tuning.}
\label{fig:finite}
\end{figure}

Table~\ref{tab:syntheticfour} also reports the terminal normalized
objective gap $F_k=[f(x_k)-f^\star]/[f(x_0)-f^\star]$, which
measures objective accuracy relative to initialization. CHB has the
smallest $D_{512}$ and $F_{512}$ in both exact-gradient examples;
the same ordering holds with scalar-half errors. These high-precision
errors describe convergence on the fixed examples, not runtime gains
or a general finite-horizon guarantee.

\begin{table}[htbp]
\centering
\caption{Terminal distances and objective gaps with exact gradients and
512 queries in 80/120-digit arithmetic on the controlled examples.
The definitions of $D_k$ and
$F_k$ give one at initialization; smaller values are better.}
\label{tab:syntheticfour}
\small
\begin{tabular}{@{}lrrrr@{}}
\toprule
& \multicolumn{2}{c}{$R=1/2$} & \multicolumn{2}{c}{$R=1/8$}\\
\cmidrule(lr){2-3}\cmidrule(l){4-5}
Method & $D_{512}$ & $F_{512}$ & $D_{512}$ & $F_{512}$\\
\midrule
CHB & $1.0945\times10^{-127}$ & $1.9181\times10^{-254}$ & $1.5152\times10^{-112}$ & $4.0340\times10^{-224}$\\
HB & $1.0954\times10^{-111}$ & $2.1334\times10^{-222}$ & $1.0954\times10^{-111}$ & $2.1334\times10^{-222}$\\
GD & $1.0432\times10^{-28}$ & $1.0882\times10^{-56}$ & $1.0432\times10^{-28}$ & $1.0882\times10^{-56}$\\
NAG-Q & $1.1209\times10^{-73}$ & $1.4786\times10^{-147}$ & $1.1209\times10^{-73}$ & $1.4786\times10^{-147}$\\
\bottomrule
\end{tabular}
\end{table}

The CHB/HB comparisons also cover $e_k=(3\varepsilon_U/2)Hx_k$
and mixing errors $e_k=(\varepsilon_U/2)Q_kHx_k$, where
$P(v_1,v_2,v_3,v_4)^T=(v_4,v_1,v_2,v_3)^T$ and $Q_k$ cycles through $P,P^T,-P$, starting with $Q_0=P$.
Above-threshold CHB/HB per-query factors are $(0.609714,0.647302)$
at $R=1/2$ and $(0.600501,0.602716)$ at $R=1/8$: both converge
above $q$. All positive synthetic budgets exceed $c_\kappa R^4$;
these selected paths do not establish a worst-case vector-error guarantee.

GD and NAG-Q calculations use mpmath 1.3.0 at 80/120 digits.
Their norm-log discrepancies are below
$1.06\times10^{-76}$; affine-state discrepancies at $k=1,2,16,64,512$
are below $4.37\times10^{-76}$. Synthetic CHB/HB objective values
exist only at saved states.

\section{An illustration with stock index returns}\label{sec:istanbul}
Use all 536 Istanbul Stock Exchange~\cite{istanbul2013} observations in
source order: seven raw international returns (SP, DAX, FTSE, NIKKEI,
BOVESPA, EU, EM) and the USD-based ISE response. Neither date nor TL-based
ISE enters the design. Before optimization, feature-only screening selected
Istanbul by the lowest raw condition number among four eligible datasets;
public candidate records document this targeted selection.
For $n=536$, design $X\in\R^{n\times7}$, response $y\in\R^n$ and
coefficients $w\in\R^7$, define the objective, Hessian and gradient by
\begin{equation}\label{eq:regressionobjective}
 f(w)=\frac{\norm{Xw-y}^2}{2n},\qquad
 H=\frac{X^TX}{n},\qquad h=\frac{X^Ty}{n},\qquad g(w)=Hw-h.
\end{equation}
Equation \eqref{eq:regressionobjective} defines contemporaneous least squares
without an intercept, centering, scaling, penalty or train/test split.

The screening eigenvalues $\widehat\lambda_i$ define enclosing bounds
$\mu=\widehat\lambda_{\min}-\tau$ and
$L=\widehat\lambda_{\max}+\tau$, where
$\tau=64u_{\rm fp}\widehat\lambda_{\max}$ and $u_{\rm fp}=2^{-52}$
is binary64 machine precision. These are floating-point safeguards,
not interval certificates; the saved high-precision spectrum also lies
inside these bounds. The resulting $\mu\approx5.23525\times10^{-6}$,
$L\approx8.56403\times10^{-4}$ and fixed
$R=0.49644090689845105$ enclose all seven eigenvalues in
\eqref{eq:class}, with six in the lower band and one in the upper.
Here $\kappa\approx163.584$. All methods use the fixed tuning rules
in Section~\ref{sec:numerics}, $w_{-1}=w_0=0$ and 512 queries;
the bounds and horizon are not retuned. Displayed endpoints are rounded;
tuning uses the full decimal values in the fixed configuration, and
$a,b$ are obtained from \eqref{eq:class}, not treated as eigenvalues.

The experiment used the earlier sufficient constant
$c_\kappa^{\rm p}=\min\{1,d_-^2/(4K_{\rm p}G_{\rm p})\}$, with
$K_{\rm p}=\sqrt{(4\kappa^2+3)^2+2}$ and
$G_{\rm p}=4\kappa(\sqrt{4\kappa^2+2}+\kappa)$.
Here the superscript ${\rm p}$ denotes the fixed protocol.
Since $\mathsf A_r,\mathsf B_r\leq2\kappa$, the constants in
\eqref{eq:auxconstants} give $c_\kappa\geq c_\kappa^{\rm p}$.
Compare exact gradients with errors $e_k=\varepsilon Q_k g(z_k)$,
where $z_k$ is each method's query point and the unchanged budget is
$\varepsilon=c_\kappa^{\rm p}R^4/2\approx1.60364486\times10^{-15}$.
For $P(v_1,\ldots,v_7)^T=(v_7,v_1,\ldots,v_6)^T$,
$Q_k$ cycles through $P,P^T,-P$ from $Q_0=P$.
Orthogonality enforces the computational error budget, which lies
in the sufficient region of Theorem~\ref{thm:main}.

Let $w^\star$ be the exact minimizer, $\widehat w$ the computed QR
reference, and $r=H\widehat w-h$ its gradient residual. For $k=0,\ldots,512$,
use the cancellation-free energy
$E_{\widehat w}(w)=\tfrac12(w-\widehat w)^TH(w-\widehat w)$ to define
\begin{equation}\label{eq:regressionmetrics}
 D_k=\frac{\norm{w_k-\widehat w}}{\norm{w_0-\widehat w}},\qquad
 F_k=\frac{E_{\widehat w}(w_k)}{E_{\widehat w}(w_0)}.
\end{equation}
The normalized distance and energy gap in \eqref{eq:regressionmetrics}
equal one initially. The exact objective difference includes the residual:
\begin{equation}\label{eq:regressiongap}
 f(w)-f(\widehat w)=E_{\widehat w}(w)+r^T(w-\widehat w).
\end{equation}
In \eqref{eq:regressiongap}, the residual term vanishes at $\widehat w=w^\star$.
Its magnitude relative to $E_{\widehat w}$ is below $4.95\times10^{-123}$
at the saved 160-digit checkpoints. The code documents reference-error
bounds and affine-state checks retaining $r$.

\begin{figure}[!htbp]
\centering
\includegraphics[width=\textwidth]{figures_v31_final/Fig3_clean.pdf}
\caption{Istanbul exact-gradient curves for \eqref{eq:regressionmetrics}.
The inset resolves the first 16 queries of the energy gap and CHB's
early overshoot. Selected mixing-error curves nearly coincide.}
\label{fig:istanbul}
\end{figure}

\begin{table}[htbp]
\centering
\caption{Istanbul QR-referenced errors from \eqref{eq:regressionmetrics}. The displayed
five significant digits coincide under both policies. Smaller is better.}
\label{tab:istanbulfour}
\small
\begin{tabular}{@{}lrrr@{}}
\toprule
Method & $D_{511}$ & $D_{512}$ & $F_{512}$\\
\midrule
CHB & $6.2203\times10^{-38}$ & $1.8162\times10^{-38}$ & $2.3430\times10^{-75}$\\
HB & $5.2508\times10^{-33}$ & $4.4980\times10^{-33}$ & $1.4922\times10^{-64}$\\
GD & $1.4418\times10^{-3}$ & $1.4243\times10^{-3}$ & $2.9875\times10^{-6}$\\
NAG-Q & $3.3370\times10^{-20}$ & $3.0419\times10^{-20}$ & $4.3029\times10^{-41}$\\
\bottomrule
\end{tabular}
\end{table}

Table~\ref{tab:istanbulfour} reports adjacent terminal queries. To assess both phases, set
$\widehat D_j=\max\{D_{2j-1},D_{2j}\}$ and
$\widehat F_j=\max\{F_{2j-1},F_{2j}\}$ for $j=1,\ldots,256$.
For exact gradients, the terminal CHB/HB ratios of
$(\widehat D_{256},\widehat F_{256})$ are
$(1.1846\times10^{-5},1.4300\times10^{-10})$; CHB/NAG-Q gives
$(1.8640\times10^{-18},5.6153\times10^{-34})$.
Under both policies, CHB has smaller distances and energy gaps than all
comparators at every recorded query 62--512.

Figure~\ref{fig:istanbul} also shows substantial early amplification:
CHB peaks at $F_k\approx62.815$ at query 5 and $D_k\approx2.9221$.
HB's peak energy gap is about 17.909; GD and NAG-Q remain at or below
initialization throughout the window.

The 16 trajectories use mpmath 1.3.0 at 120/160 digits, with 160 primary.
Norm-log discrepancies are below $1.10\times10^{-82}$, affine-state errors
below $1.57\times10^{-83}$ and normalized QR residuals below
$2.29\times10^{-122}$. One selected Hessian and a tiny error budget
establish neither practical noise tolerance nor wall-clock gains.

\section{Interpretation and limitations}\label{sec:discussion}
The constants in Theorem~\ref{thm:main} are conservative;
\eqref{eq:upperlimit} concerns only the scalar obstruction. Exact
vector-error constants and the full rate--robustness tradeoff under
noise-dependent tuning remain open. The low-dimensional examples
illustrate the analysis and transients; they do not assess scalability.

At fixed $\kappa$ within the theorem's error range, $\delta$ has order
$R^2$ and $\gamma$ stays away from zero as $R\downarrow0$, so
the prefactor in \eqref{eq:prefactor} grows as $R^{-2}$.
This bound guarantees neither a short-horizon advantage nor monotone
objective decrease. Extending the fixed-Hessian, equal-width-band
analysis to unequal bands, nonquadratic objectives, stochastic errors,
or explicit crossover horizons remains open.

\backmatter
\section*{Statements and Declarations}
\subsection*{Funding}
No funding was received for this work.
\subsection*{Competing interests}
The authors declare no competing interests.
\subsection*{Data and code availability}
{\raggedright
The Istanbul Stock Exchange data are available from UCI~\cite{istanbul2013}
under CC BY 4.0. Code, fixed configurations, candidate-selection records,
source data and accepted results are public at
\url{https://github.com/byuzbasi/periodic-momentum}.
Reproduction guides and figures are available at
\url{https://byuzbasi.github.io/periodic-momentum/}.\par}
\subsection*{Use of generative AI}
During manuscript preparation, ChatGPT
(OpenAI) was used for English-language editing and stylistic refinement.
The author reviewed all content and takes full responsibility for the
accuracy and scientific integrity of the manuscript.

\begin{thebibliography}{12}
\raggedright
\bibitem{goujaud2022}
Goujaud, B., Scieur, D., Dieuleveut, A., Taylor, A.B., Pedregosa, F.:
Super-acceleration with cyclical step-sizes.
In: Proceedings of the 25th International Conference on Artificial
Intelligence and Statistics, PMLR \textbf{151}, 3028--3065 (2022).
\url{https://proceedings.mlr.press/v151/goujaud22a.html}

\bibitem{oymak2021}
Oymak, S.: Provable super-convergence with a large cyclical learning rate.
IEEE Signal Process. Lett. \textbf{28}, 1645--1649 (2021).
\url{https://doi.org/10.1109/LSP.2021.3101131}

\bibitem{agarwal2021}
Agarwal, N., Goel, S., Zhang, C.: Acceleration via fractal learning rate schedules.
In: Proceedings of the 38th International Conference on Machine Learning,
PMLR \textbf{139}, 87--99 (2021).
\url{https://proceedings.mlr.press/v139/agarwal21a.html}

\bibitem{gannot2022}
Gannot, O.: A frequency-domain analysis of inexact gradient methods.
Math. Program. \textbf{194}, 975--1016 (2022).
\url{https://doi.org/10.1007/s10107-021-01665-8}

\bibitem{cyrus2018}
Cyrus, S., Hu, B., Van Scoy, B., Lessard, L.:
A robust accelerated optimization algorithm for strongly convex functions.
In: 2018 Annual American Control Conference, pp. 1376--1381 (2018).
\url{https://doi.org/10.23919/ACC.2018.8430824}

\bibitem{gurbuz2025}
G\"urb\"uzbalaban, M.: Robustly stable accelerated momentum methods with a
near-optimal $L_2$ gain and $H_\infty$ performance.
Math. Oper. Res., published online 5 December 2025.
\url{https://doi.org/10.1287/moor.2023.0321}

\bibitem{khanh2026}
Khanh, P.D., Mordukhovich, B.S., Tran, D.B.:
Convergence of first-order algorithms with momentum from the perspective
of an inexact gradient descent method.
J. Optim. Theory Appl. \textbf{208}, article 40 (2026).
\url{https://doi.org/10.1007/s10957-025-02864-5}

\bibitem{miller2026}
Miller, J., Jakob, F., Scherer, C., Iannelli, A.:
Analysis and synthesis of switched optimization algorithms.
Nonlinear Anal. Hybrid Syst. \textbf{62}, article 101788 (2026).
\url{https://doi.org/10.1016/j.nahs.2026.101788}

\bibitem{orvieto2025}
Orvieto, A.: An accelerated Lyapunov function for Polyak's heavy-ball
on convex quadratics. Optim. Lett. \textbf{19}, 307--328 (2025).
\url{https://doi.org/10.1007/s11590-024-02119-8}

\bibitem{deklerk2017}
de Klerk, E., Glineur, F., Taylor, A.B.:
On the worst-case complexity of the gradient method with exact line search
for smooth strongly convex functions.
Optim. Lett. \textbf{11}, 1185--1199 (2017).
\url{https://doi.org/10.1007/s11590-016-1087-4}
\bibitem{lessard2016}
Lessard, L., Recht, B., Packard, A.:
Analysis and design of optimization algorithms via integral quadratic constraints.
SIAM J. Optim. \textbf{26}(1), 57--95 (2016).
\url{https://doi.org/10.1137/15M1009597}

\bibitem{istanbul2013}
Akbilgic, O.: Istanbul Stock Exchange [Dataset].
UCI Machine Learning Repository (2013).
\url{https://doi.org/10.24432/C54P4J}

\end{thebibliography}
\end{document}
